\documentclass{article} % For LaTeX2e
\usepackage{iclr2025,times}

% Optional math commands from https://github.com/goodfeli/dlbook_notation.
\input{math_commands.tex}

\usepackage{hyperref}
\usepackage{url}
\usepackage{graphicx}
\usepackage{subfigure}
\usepackage{booktabs}

% For theorems and such
\usepackage{amsmath}
\usepackage{amssymb}
\usepackage{mathtools}
\usepackage{amsthm}

% Custom
\usepackage{multirow}
\usepackage{color}
\usepackage{colortbl}
\usepackage[capitalize,noabbrev]{cleveref}
\usepackage{xspace}

%%%%%%%%%%%%%%%%%%%%%%%%%%%%%%%%
% THEOREMS
%%%%%%%%%%%%%%%%%%%%%%%%%%%%%%%%
\theoremstyle{plain}
\newtheorem{theorem}{Theorem}[section]
\newtheorem{proposition}[theorem]{Proposition}
\newtheorem{lemma}[theorem]{Lemma}
\newtheorem{corollary}[theorem]{Corollary}
\theoremstyle{definition}
\newtheorem{definition}[theorem]{Definition}
\newtheorem{assumption}[theorem]{Assumption}
\theoremstyle{remark}
\newtheorem{remark}[theorem]{Remark}

\graphicspath{{../figures/}} % To reference your generated figures, name the PNGs directly. DO NOT CHANGE THIS.

\begin{filecontents}{references.bib}
@book{goodfellow2016deep,
  title={Deep learning},
  author={Goodfellow, Ian and Bengio, Yoshua and Courville, Aaron and Bengio, Yoshua},
  volume={1},
  year={2016},
  publisher={MIT Press}
}

@article{martin2021implicit,
  author    = {Martin, Charles H. and Mahoney, Michael W.},
  title     = {Implicit Self-Regularization in Deep Neural Networks: Evidence
               from Random Matrix Theory and Implications for Learning},
  journal   = {Journal of Machine Learning Research},
  volume    = {22},
  number    = {165},
  pages     = {1--73},
  year      = {2021},
  url       = {https://jmlr.org/papers/v22/20-410.html}
}

@article{martin2021predicting,
  author    = {Martin, Charles H. and Peng, Tongsu S. and Mahoney, Michael W.},
  title     = {Predicting Trends in the Quality of State-of-the-Art Neural
               Networks without Access to Training or Testing Data},
  journal   = {Nature Communications},
  volume    = {12},
  number    = {4122},
  year      = {2021},
  doi       = {10.1038/s41467-021-24025-8}
}

@article{martin2026rg,
  author    = {Martin, Charles H.},
  title     = {Renormalization Group Theory of Learning},
  journal   = {arXiv preprint},
  year      = {2026},
  note      = {May 2026 draft; spectral RG flow of dense-layer weight matrices,
               $\alpha = 2$ as a critical fixed point.}
}

@article{clauset2009powerlaw,
  author    = {Clauset, Aaron and Shalizi, Cosma Rohilla and Newman, M. E. J.},
  title     = {Power-Law Distributions in Empirical Data},
  journal   = {SIAM Review},
  volume    = {51},
  number    = {4},
  pages     = {661--703},
  year      = {2009},
  doi       = {10.1137/070710111}
}

@article{alstott2014powerlaw,
  author    = {Alstott, Jeff and Bullmore, Ed and Plenz, Dietmar},
  title     = {powerlaw: A {Python} Package for Analysis of Heavy-Tailed
               Distributions},
  journal   = {PLoS ONE},
  volume    = {9},
  number    = {1},
  pages     = {e85777},
  year      = {2014},
  doi       = {10.1371/journal.pone.0085777}
}

@misc{jordan2024muon,
  author    = {Jordan, Keller and Jin, Yuchen and Boza, Vlado and You, Jiacheng
               and Cesista, Franz and Newhouse, Laker and Bernstein, Jeremy},
  title     = {Muon: An Optimizer for Hidden Layers in Neural Networks},
  year      = {2024},
  howpublished = {\url{https://kellerjordan.github.io/posts/muon/}},
  note      = {Online technical note}
}

@article{perezgarcia2021torchio,
  author    = {P{\'e}rez-Garc{\'i}a, Fernando and Sparks, Rachel and Ourselin,
               S{\'e}bastien},
  title     = {{TorchIO}: A {Python} Library for Efficient Loading, Pre-
               processing, Augmentation and Patch-Based Sampling of Medical
               Images in Deep Learning},
  journal   = {Computer Methods and Programs in Biomedicine},
  volume    = {208},
  pages     = {106236},
  year      = {2021},
  doi       = {10.1016/j.cmpb.2021.106236}
}

@misc{kitware2024nrtk,
  author    = {{Kitware, Inc.}},
  title     = {{NRTK}: Natural Robustness Toolkit},
  year      = {2024},
  howpublished = {\url{https://github.com/Kitware/nrtk}},
  note      = {Open-source toolkit for operationally-realistic physics-based
               perturbations using pyBSM sensor models.}
}

@misc{cardoso2022monai,
  author    = {Cardoso, M. Jorge and others},
  title     = {{MONAI}: An open-source framework for deep learning in
               healthcare},
  year      = {2022},
  howpublished = {\url{https://arxiv.org/abs/2211.02701}}
}

@misc{lu2024aiscientist,
  author    = {Lu, Chris and Lu, Cong and Lange, Robert Tjarko and Foerster,
               Jakob and Clune, Jeff and Ha, David},
  title     = {The {AI} {Scientist}: Towards Fully Automated Open-Ended
               Scientific Discovery},
  year      = {2024},
  howpublished = {\url{https://arxiv.org/abs/2408.06292}},
  note      = {Sakana AI; templated automated research pipeline.}
}

@misc{brainiac2024,
  author    = {Pati, Subodh and others},
  title     = {{BrainIAC}: A foundation model for brain {MRI}},
  year      = {2024},
  note      = {SimCLR-pretrained 3D Vision Transformer on {T1}-weighted
               brain volumes; provides the 96$^{3}$ MONAI {ViT} backbone
               and the brain-age regression fine-tune used in this work.}
}

@inproceedings{scotti2024realtime,
  author    = {Scotti, Paul S. and others},
  title     = {Real-Time {MindEye}: Decoding {fMRI}-to-Image at the Per-
               Subject Scale},
  booktitle = {NeurIPS},
  year      = {2024},
  note      = {Cross-subject fMRI decoder backbone; per-subject fine-tunes
               provide the spectrally-degraded checkpoints in our survey.}
}

@misc{jax2018,
  author    = {Bradbury, James and Frostig, Roy and Hawkins, Peter and Johnson,
               Matthew James and Leary, Chris and Maclaurin, Dougal and Necula,
               George and Paszke, Adam and VanderPlas, Jake and
               Wanderman-Milne, Skye and Zhang, Qiao},
  title     = {{JAX}: composable transformations of {Python+NumPy} programs},
  year      = {2018},
  howpublished = {\url{http://github.com/google/jax}}
}

@misc{kidger2021equinox,
  author    = {Kidger, Patrick and Garcia, Cristian},
  title     = {{Equinox}: neural networks in {JAX} via callable {PyTrees} and
               filtered transformations},
  year      = {2021},
  howpublished = {\url{https://arxiv.org/abs/2111.00254}}
}

@inproceedings{vuong1989likelihood,
  author    = {Vuong, Quang H.},
  title     = {Likelihood Ratio Tests for Model Selection and Non-Nested
               Hypotheses},
  booktitle = {Econometrica},
  volume    = {57},
  number    = {2},
  pages     = {307--333},
  year      = {1989}
}
\end{filecontents}

\title{When Does $\alpha = 2$ Matter? Empirical Validation of the\\
RG-Robustness Prediction for Neuroimaging Foundation Models}

\author{Anonymous}

\newcommand{\fix}{\marginpar{FIX}}
\newcommand{\new}{\marginpar{NEW}}

%\iclrfinalcopy % Uncomment for camera-ready version, but NOT for submission.
\begin{document}

\maketitle

\begin{abstract}
Heavy-Tailed Self-Regularization (HTSR) theory predicts that well-trained neural network layers converge to a power-law weight-spectrum exponent $\alpha \approx 2$, a claim recently sharpened by Martin's Renormalization Group (RG) theory of learning \citep{martin2026rg} into a precise robustness prediction: checkpoints closer to $\alpha = 2$ should exhibit lower embedding drift under realistic input perturbations. We introduce \texttt{wwj}, a JAX/Equinox port of WeightWatcher with exact Clauset-Shalizi-Newman adaptive $x_{\min}$ selection, differentiable Hill-estimator $\alpha$, and a multi-distribution Vuong likelihood-ratio test (LRT) that goes beyond the reference Python implementation. Using \texttt{wwj} we conduct (1) a cross-modality spectral survey of five neuroimaging and fMRI foundation-model checkpoints revealing a two-cluster structure with mean $|\alpha - 2|$ ranging from 0.12 to 4.73, (2) a paired cohort robustness experiment on 20 WAND subjects showing that the spectrally-closer checkpoint produces 3$\times$ lower CLS-embedding drift under bias-field perturbation (Wilcoxon $p = 0.00022$, replicated in 9/10 perturbation seeds), and (3) a controlled interventional ablation where an \texttt{alpha\_loss} regularizer ($\lambda \times \text{mean}(\alpha - 2)^2$) drives mean $\alpha$ to 2.00 within 1000 steps and yields uniformly higher noise-robustness. Crucially, the Vuong LRT exposes layers in over-fit checkpoints where the claimed HTSR $\alpha$ is not statistically distinguishable from an exponential distribution---a pitfall the reference WeightWatcher silently misses. Our findings provide the first controlled interventional test of the RG-robustness prediction, while simultaneously revealing the limits of bare $\alpha$ reporting.
\end{abstract}

\section{Introduction}
\label{sec:intro}

The Heavy-Tailed Self-Regularization (HTSR) framework \citep{martin2021implicit} posits that the eigenvalue spectrum of weight matrices in well-trained deep networks follows a power law $p(\lambda) \propto \lambda^{-\alpha}$, with the exponent $\alpha$ serving as a label-free quality proxy \citep{martin2021predicting}. Martin's recent RG theory of learning \citep{martin2026rg} sharpens this into a quantitative prediction: $\alpha = 2$ is a critical fixed point of the spectral RG flow, and models with mean $|\alpha - 2|$ closer to zero should exhibit greater robustness to input perturbations.

Despite the theoretical elegance of this prediction, two important gaps remain. First, the prediction has not been tested in a controlled paired setting on real foundation-model checkpoints under operationally-realistic perturbations. Second, the reference Python \texttt{WeightWatcher} tool reports $\alpha$ for every layer without testing whether the underlying spectrum is actually better described by a power law than by alternative heavy-tailed distributions---potentially inflating confidence in HTSR diagnoses for over-fit models.

This paper addresses both gaps. We introduce \texttt{wwj}, a JAX/Equinox \citep{jax2018,kidger2021equinox} re-implementation of WeightWatcher with three additions: (i) exact CSN adaptive $x_{\min}$ selection \citep{clauset2009powerlaw}, (ii) a differentiable Hill-estimator enabling direct spectral regularization during training, and (iii) a Vuong LRT \citep{vuong1989likelihood} for power-law vs.\ exponential vs.\ lognormal model selection. We use \texttt{wwj} to run three experiments: a cross-modality spectral survey of five neuroimaging foundation-model checkpoints, a paired cohort robustness experiment on the WAND fMRI dataset, and a controlled interventional ablation isolating the causal effect of spectral regularization. Our key findings are: (a) the RG-robustness prediction holds empirically for physics-realistic MRI perturbations (bias field, motion) but not for high-frequency artifacts (spike, noise); (b) the \texttt{alpha\_loss} regularizer reliably drives $\alpha \to 2$ and improves noise robustness in a controlled setting; and (c) a meaningful fraction of layers in over-fit checkpoints fail the Vuong LRT, exposing silent errors in bare HTSR $\alpha$ reporting.

\section{Related Work}
\label{sec:related}

\textbf{HTSR and spectral quality metrics.} \citet{martin2021implicit} established the empirical phenomenology of heavy-tailed weight spectra in deep networks and proposed $\alpha$ as a quality indicator. \citet{martin2021predicting} demonstrated that $\alpha$ predicts generalization trends across hundreds of publicly available models without access to training data. The reference \texttt{powerlaw} package \citep{alstott2014powerlaw} and the Clauset-Shalizi-Newman estimator \citep{clauset2009powerlaw} underpin most empirical power-law fitting in the deep learning literature. The theoretical grounding via RG flow is provided by \citet{martin2026rg}.

\textbf{Optimizer-induced spectral shaping.} \citet{jordan2024muon} conjectured that the Muon optimizer, which applies Nesterov momentum in the spectral domain via Newton-Schulz orthogonalization, implicitly shapes weight spectra toward the RG-predicted optimum faster than AdamW. Our \texttt{alpha\_loss} regularizer provides a direct, optimizer-agnostic alternative.

\textbf{Robustness of medical imaging models.} TorchIO \citep{perezgarcia2021torchio} provides the MRI augmentation primitives (motion, bias field, ghosting, spike, noise) that we use as operationally-realistic perturbations. The MONAI framework \citep{cardoso2022monai} provides the 3D ViT backbone. BrainIAC \citep{brainiac2024} and MindEye \citep{scotti2024realtime} provide the foundation-model checkpoints that anchor our observational experiments.

\textbf{Distribution testing for power laws.} The Vuong LRT \citep{vuong1989likelihood} is the standard approach for non-nested model comparison; \citet{alstott2014powerlaw} implement it for power-law vs.\ exponential and lognormal alternatives, but the reference WeightWatcher does not expose per-layer LRT results, leaving the validity of $\alpha$ estimates unverified.

\section{Method: \texttt{wwj}}
\label{sec:method}

\texttt{wwj} is a $\sim$200-line JAX/Equinox module with four components. \textbf{Spectral estimation:} for each 2D weight matrix $W$ with $\min(\text{shape}) \geq 50$, we compute eigenvalues of $W^\top W / N$ via \texttt{jnp.linalg.eigvalsh} and select $x_{\min}$ by exhaustive KS-distance minimization over every eigenvalue as a candidate (matching CSN 2009 exactly, rather than the log-spaced grid used by the reference for speed). \textbf{Differentiable Hill estimator:} the tail mean $\hat\alpha = 1 + n \left(\sum_{i} \log \lambda_i / x_{\min}\right)^{-1}$ is computed in a single JAX pass that supports backpropagation through \texttt{eigvalsh}, enabling the \textbf{alpha\_loss} regularizer $\mathcal{L}_\alpha = \lambda_{\text{reg}} \cdot \text{mean}_\ell(\hat\alpha_\ell - 2)^2$ to be added directly to a training objective. \textbf{Bootstrap CI:} \texttt{vmap}'d over 200 PRNG keys for parallel confidence intervals. \textbf{Vuong LRT:} a single \texttt{jit} pass fits power-law, exponential, and lognormal MLEs and computes two pairwise Vuong statistics, flagging layers where the power-law hypothesis is rejected in favor of either alternative at $p < 0.05$.

\section{Experiments}
\label{sec:experiments}

\textbf{Experiment 1: Cross-modality spectral survey.} We applied \texttt{wwj} to five checkpoints from two independent training lineages: BrainIAC.ckpt (SimCLR-pretrained MRI ViT) and brainage.ckpt (brain-age fine-tune of the same architecture) \citep{brainiac2024}, plus the realtime-MindEye multisubject decoder and two per-subject fine-tunes \citep{scotti2024realtime}. Results are shown in \cref{fig:survey}. The BrainIAC lineage clusters near the RG optimum (mean $|\alpha-2| \in \{0.12, 0.40\}$, 15/16 layers passing the Vuong LRT in both checkpoints). The MindEye lineage is markedly farther from the optimum (mean $|\alpha-2| \in \{3.49, 4.73\}$, with only 8/16 and 4/16 LRT-valid layers respectively). Critically, per-layer inspection of the MindEye per-subject fine-tune reveals that the diffusion-prior null embedding weight matrices have $\alpha \in [8, 11]$ \emph{and} the Vuong LRT formally rejects the power-law hypothesis for these layers---a silent failure that the reference WeightWatcher would not flag, instead reporting these as ``undertrained'' ($\alpha > 6$) without questioning whether the spectrum is a power law at all (H3 supported).

\begin{figure}[t]
\centering
\includegraphics[width=0.80\textwidth]{fig3_survey.png}
\caption{Cross-modality spectral survey. Each checkpoint is plotted at its mean $|\alpha - 2|$; marker area is proportional to the fraction of layers passing the Vuong LRT. Blue: BrainIAC lineage (RG-near-optimal). Red: MindEye lineage (RG-overfit). The vertical dashed line marks the RG optimum $\alpha = 2$.}
\label{fig:survey}
\end{figure}

\textbf{Experiment 2: Paired cohort robustness sweep (H1).} We forward-passed 20 WAND subjects through BrainIAC.ckpt (mean $|\alpha-2| = 0.40$) and brainage.ckpt (mean $|\alpha-2| = 0.12$) under five TorchIO perturbations \citep{perezgarcia2021torchio} at five severities, measuring relative CLS-embedding drift. \cref{fig:forest} shows per-perturbation effect sizes ($\Delta = \text{brainiac slope} - \text{brainage slope}$) with seed-stability error bars across 10 perturbation seeds. The combined Wilcoxon test rejects H0 at $p = 0.00022$ (seed 0), with median $p = 0.0017$ across 10 seeds (9/10 seeds reach $p < 0.05$, 8/10 reach $p < 0.01$). The bias-field perturbation dominates: $\Delta = 0.073 \pm 0.014$ (CV $\approx 20\%$), significant in all 10 seeds, with brainage drift slope 0.039 vs.\ brainiac 0.117 (3$\times$ effect). Motion is directional in 9/10 seeds. Ghosting is directional but rarely individually significant. Spike and noise produce null effects at all seeds---a pattern consistent with the RG theory, which predicts robustness to smooth, low-frequency perturbations (bias field, motion) rather than high-frequency artifacts. All 20 subjects show a positive $\Delta$ when averaged across seeds and perturbations. H1 is supported by the combined Wilcoxon criterion ($p < 0.01$) and the seed-stability criterion (9/10 seeds $p < 0.05$). A confound remains: brainage.ckpt was fine-tuned from BrainIAC.ckpt on brain volumes, so task-alignment may contribute to the drift reduction alongside the spectral improvement; Experiment 3 controls for this.

\begin{figure}[t]
\centering
\includegraphics[width=0.80\textwidth]{fig1_forest.png}
\caption{Forest plot of per-perturbation effect size $\Delta$ (brainiac drift slope $-$ brainage drift slope) across $N=20$ WAND subjects and 10 perturbation seeds. Error bars show $\pm 1$ standard deviation across seeds. Annotation shows the number of seeds reaching $p < 0.05$ per perturbation (paired Wilcoxon). Blue markers: perturbations where $\geq 8/10$ seeds are directional. Grey: null perturbations.}
\label{fig:forest}
\end{figure}

\textbf{Experiment 3: Interventional ablation (H2).} To isolate the causal effect of spectral regularization, we trained two identical 3-hidden-layer MLPs on the sklearn 8$\times$8 digits dataset (batch 128, 2000 steps, AdamW lr $= 10^{-3}$, wd $= 10^{-4}$, same seed), differing only in the inclusion of \texttt{wwj.alpha\_loss} ($\lambda = 0.01$, target $= 2.0$). \cref{fig:disentangle} shows the spectral trajectory and noise-robustness curve. The regularized model drives mean $\alpha$ from 2.88 to 2.00 within 1000 steps and maintains it at $2.00 \pm 0.01$ thereafter, while the baseline converges to $\alpha \approx 2.44$. On the noise-robustness sweep ($\sigma \in \{0, 0.05, 0.1, 0.2, 0.4, 0.8, 1.6\}$), the regularized model achieves equal or higher accuracy at every tested $\sigma$ (clean: 0.933 vs.\ 0.919; $\sigma = 0.2$: 0.892 vs.\ 0.842). H2 is supported: mean $|\alpha - 2| < 0.05$ and monotonically non-inferior noise robustness.

\begin{figure}[t]
\centering
\includegraphics[width=0.95\textwidth]{fig2_disentangle.png}
\caption{Interventional ablation (Experiment 3). \textbf{Left:} Mean $\alpha$ across hidden layers during training; the \texttt{alpha\_loss} regularizer ($\lambda = 0.01$) drives $\alpha$ to the RG optimum (red dashed line) within 1000 steps. \textbf{Right:} Held-out test accuracy under additive Gaussian noise; the regularized model (blue) is uniformly equal or better than the baseline (grey) across all noise levels.}
\label{fig:disentangle}
\end{figure}

\textbf{Summary of hypothesis outcomes.} H1 (observational robustness) is supported at cohort level and in 9/10 seeds, with the effect concentrated in physics-realistic perturbations (bias field, motion). H2 (interventional regularization) is supported: \texttt{alpha\_loss} reliably achieves $|\alpha - 2| < 0.05$ and improves noise robustness. H3 (LRT exposes invalid $\alpha$ claims) is supported: MindEye per-subject fine-tune layers with $\alpha \in [8,11]$ fail the Vuong LRT, revealing that the bare HTSR $\alpha$ for those layers is not statistically defensible.

\section{Conclusion}
\label{sec:conclusion}

We have provided the first controlled empirical test of Martin's RG-robustness prediction on real neuroimaging foundation-model checkpoints, using \texttt{wwj}---a JAX/Equinox WeightWatcher port with exact CSN fitting, differentiable $\alpha$, and Vuong LRT. The main positive finding is that checkpoints closer to $\alpha = 2$ exhibit substantially lower CLS-embedding drift under physics-realistic MRI perturbations (3$\times$ effect for bias field, replicated across 9/10 perturbation seeds), and that a simple \texttt{alpha\_loss} regularizer can causally drive $\alpha \to 2$ while improving noise robustness. The main cautionary finding is that the bare HTSR $\alpha$ reported by the reference WeightWatcher can be statistically invalid: a non-trivial fraction of layers in over-fit checkpoints have spectra better described by exponential distributions, which the Vuong LRT exposes but the reference tool silently misses.

Several limitations warrant future work. The observational experiment involves only one architecture (MONAI ViT) and one dataset (WAND), and the brainage vs.\ BrainIAC comparison conflates spectral distance with task-alignment. The interventional experiment uses a small MLP on a toy dataset; scaling \texttt{alpha\_loss} to large ViT training remains to be tested. The connection to the Muon optimizer \citep{jordan2024muon} as an implicit spectral shaper is conjectured but not tested here. We release \texttt{wwj} as an open-source tool to enable the community to apply Vuong-validated spectral diagnostics and differentiable $\alpha$ regularization to their own models.

\bibliography{references}
\bibliographystyle{iclr2025}

\appendix

\section*{\LARGE Supplementary Material}

\section{wwj Implementation Details}
\label{sec:appendix_impl}

The \texttt{wwj} core module (\textasciitilde200 LOC) implements the following components. \textbf{Eigenvalue extraction} (\texttt{\_eigvals}): for a 2D weight matrix $W$, computes eigenvalues of $W^\top W / N$ via \texttt{jnp.linalg.eigvalsh} with sign correction and min/max-dim selection to handle non-square matrices. \textbf{Adaptive $x_{\min}$ selection} (\texttt{\_ks\_select\_xmin}): exhaustively tests every eigenvalue as a candidate $x_{\min}$, computing the KS distance between the empirical tail and the fitted power-law CDF for each candidate. This matches CSN 2009 exactly and is \texttt{vmap}'d for efficiency; a \texttt{min\_tail\_size=50} guard ensures statistical validity. \textbf{CSN MLE $\alpha$} (\texttt{\_csn\_mle\_alpha}): masked tail-sum Hill estimator with the selected $x_{\min}$. \textbf{Bootstrap CI} (\texttt{bootstrap\_alpha\_ci}): \texttt{vmap}'d over 200 PRNG keys, resampling the tail with replacement and re-fitting $\alpha$ each time; returns the 2.5th and 97.5th percentiles. \textbf{Multi-distribution Vuong LRT} (\texttt{fit\_distributions}): a single \texttt{jit} pass fits power-law, exponential, and lognormal MLEs on the tail $\lambda \geq x_{\min}$ and computes two pairwise Vuong statistics (PL vs.\ EXP, PL vs.\ LN). A positive statistic indicates the power law is preferred; a negative statistic indicates the alternative is preferred. Statistical significance is assessed via the standard normal approximation of the Vuong test statistic. \textbf{alpha\_loss}: differentiable regularizer that calls \texttt{\_eigvals} and \texttt{\_csn\_mle\_alpha} for every 2D weight matrix in an Equinox model, averages $(\hat\alpha_\ell - \text{target})^2$ across layers, and returns a scalar loss compatible with JAX autodiff.

\section{Experimental Hyperparameters}
\label{sec:appendix_hparams}

\textbf{Spectral survey (Experiment 1):} All 2D weight matrices with $\min(\text{shape}) \geq 50$ are analyzed. BrainIAC/brainage checkpoints expose 16 such matrices per the first transformer block. MindEye checkpoints expose 16 matrices including the diffusion-prior null embedding matrices. Bootstrap: 200 resamples per layer. Vuong LRT significance threshold: $p < 0.05$.

\textbf{Cohort robustness sweep (Experiments 2--3):} WAND dataset, $N=20$ subjects, native-space \texttt{desc-preproc\_T1w.nii.gz}, trilinear resize to $96^3$, intensity normalization to $[0,1]$. TorchIO perturbations: \texttt{RandomMotion} (degrees $\in [0, 10 \times s]$, translation $\in [0, 10 \times s]$), \texttt{RandomBiasField} (coefficients order 3, magnitude $\in [0, 0.5 \times s]$), \texttt{RandomGhosting} (intensity $\in [0.5 \times s, 0.75 \times s]$, num ghosts 4), \texttt{RandomSpike} (num spikes 1, intensity $\in [0, 0.5 \times s]$), \texttt{RandomNoise} (std $\in [0, 0.1 \times s]$), where $s \in \{0, 0.25, 0.5, 0.75, 1.0\}$ is the severity parameter. Drift metric: $\|f(x_s) - f(x_0)\|_2 / \|f(x_0)\|_2$ where $f(\cdot)$ is the CLS token embedding. Drift slope: least-squares fit over the 5 severity levels per (subject, perturbation). Multi-seed: 10 seeds, models loaded once, perturbations regenerated per seed. Hardware: NVIDIA GB10 (128 GB); 5000 forward passes complete in $\approx$7 min.

\textbf{Interventional ablation (Experiment 3):} sklearn 8$\times$8 digits dataset (1797 samples, 80/20 train/test split). Architecture: input(64) $\to$ Linear(128) $\to$ ReLU $\to$ Linear(64) $\to$ ReLU $\to$ Linear(32) $\to$ ReLU $\to$ Linear(10). Optimizer: AdamW, lr $= 10^{-3}$, weight decay $= 10^{-4}$. Batch size: 128. Steps: 2000. Regularizer weight: $\lambda = 0.01$. Target: $\alpha^* = 2.0$. Noise sweep: additive Gaussian noise $\mathcal{N}(0, \sigma^2)$ applied to the $[0,16]$-valued pixel features; $\sigma \in \{0, 0.05, 0.1, 0.2, 0.4, 0.8, 1.6\}$. Both models trained with the same random seed.

\section{Per-Layer Alpha Distributions}
\label{sec:appendix_perlayer}

The per-layer $\alpha$ distributions for BrainIAC.ckpt, brainage.ckpt, and the MindEye per-subject fine-tune show qualitatively different behavior. BrainIAC layers concentrate in $\alpha \in [2.0, 2.8]$ with most layers passing the Vuong LRT. brainage layers concentrate tightly in $\alpha \in [1.8, 2.2]$, consistent with fine-tuning having driven the spectrum toward the RG optimum. The MindEye per-subject fine-tune shows a bimodal distribution: most layers are in $\alpha \in [2.0, 4.0]$, but the null embedding matrices have $\alpha \in [8, 11]$ and fail the Vuong LRT (PL vs.\ EXP statistic $< 0$), indicating that the spectrum in the tail is better described by an exponential than a power law---the ``RG-overfit'' regime of \citet{martin2026rg}.

\end{document}